\section{The arithmetic coefficients of the doubled kernel}

% BEGIN O74E1
\paragraph{O74E1. The graph indexing set.}
Let $X$ be the source Hermitian datum over the CM field $E$, with
nonsingular Hermitian form $h$ on $E^r$, Gram matrix $G$, and
polarization parameter $\kappa_0\ne0$. The indices of its doubled
Eisenstein series are the bottom row spaces $U$ of the rational
split unitary matrices $\sigma$ of size $2r$. Put
\[
 T=\begin{pmatrix}1&(\kappa_0/4)G\\-1&(\kappa_0/4)G\end{pmatrix},
 \qquad i_\gamma=T^{-1}\operatorname{diag}
       (1,(\gamma^{-1})^{t})T .
\]
For every $U$, exactly one of the following holds. Either $U$ is the
bottom row space of $i_\gamma$ for a unique unitary matrix $\gamma$ with entries in $E$,
or $U$ contains $(\xi,0)T$ for some nonzero $\xi$ with $h(\xi,\xi)=0$.

Indeed, $T$ is invertible. Multiplication in blocks gives the pairing
on the inverse image $I=UT^{-1}$ as $(\kappa_0/2)(h\oplus(-h))$.
The bottom rows of $\sigma$ are independent and mutually isotropic,
so $I$ is isotropic of dimension $r$. If projection to the second
summand has a nonzero kernel, any nonzero $(\xi,0)$ in it works.
Otherwise this projection is an isomorphism and $I=\{(Dv,v)\}$.
Isotropy gives $h(Dv,Dw)=h(v,w)$ for all $v,w$. Nondegeneracy makes
$D$ injective and hence a unitary automorphism. To identify the
index with the source matrix, multiply a row $(0,c)$ by $i_\gamma$.
The result is $(u,\gamma^{-1}u)T$, where
$u=c((\kappa_0/2)G)^{-1}$ ranges through all rows and
$\gamma^{-1}$ acts on the coordinate vector as a column.
Choosing $\gamma=D$ gives the claimed space and recovers $\gamma$
uniquely. Its projection to the second summand is injective.
% END O74E1

% BEGIN O74E2
\paragraph{O74E2. A common nonzero scalar at the origin.}
Fix the symplectic basis of $X.\Lambda$, the labeling
$\{1,\ldots,g\}\simeq\Psi\times\{1,\ldots,r\}$ and even $n$.
Write $D(\phi,\psi)$ for the native doubled test function.
Use its theta product formula, sesquilinearity and Heisenberg
intertwining as inputs here. There is a single $a\in\mathbb C^\times$
for all rational Schwartz--Bruhat tests such that
\[
 B(\phi,\psi):=D(\phi,\psi)(0)=a\langle\phi,\psi\rangle,
 \qquad \langle\phi,\psi\rangle=\int\phi\overline\psi .
\]
In rational symplectic coordinates put $d=ng$. The simultaneous action
of a rational Heisenberg displacement on both arguments becomes a pure
modulation on $D$: the translation block is $y-y=0$. Its value at zero
is one as a multiplier. Every translation and modulation on
$\mathbb Q^d$ occurs because the chosen basis identifies the rational
phase space with $\mathbb Q^{2d}$. Thus $B$ is invariant under both
simultaneous actions.

At a common positive level $R$, expand the tests in the indicators
$e_u$ of the cosets $u+R\mathbb Z^d$ inside $R^{-1}\mathbb Z^d$.
On $e_u$, modulation by the $j$th coordinate divided by $R$ acts by
$\exp(2\pi i u_j/R)$. For distinct cosets some $j$ distinguishes the
two characters. Invariance and sesquilinearity give $B(e_u,e_v)=0$
for $u\ne v$. Translation makes all diagonal values equal. The
source Haar normalization gives mass $R^{-d}$ to each such coset,
so the proportionality scalar at this level is $R^d B(e_0,e_0)$.
The test $u_0=\mathbf1_{\mathbb Z^d}$ has norm squared one and belongs
to every level, hence this scalar always equals $B(u_0,u_0)$.
Any two tests have a common level by taking a multiple of their levels.

To prove $a\ne0$, first $D(u_0,u_0)\ne0$. Otherwise evaluate the
theta product identity on equal frames with the second elliptic
argument the conjugate of the first. It gives the absolute square
of $\Theta_{u_0}$ identically zero. At any regular frame the period
change is invertible, the Gaussian prefactor is nonzero, and the
Fourier coefficient at frequency zero of the ordinary theta series
is $u_0(0)=1$. This is impossible. If $a=0$, apply $B$ to the
first test translated by the native displacement $-t$ and to $u_0$
in the second slot. Doubling sends this displacement to the Hermitian
Heisenberg operator with translation $-t$. Its value at zero is
$D(u_0,u_0)(t)$ times a nonzero phase for every
$t\in\operatorname{Mat}_{n,r}(E)$. The assumed vanishing would give
$D(u_0,u_0)=0$, a contradiction.
% END O74E2

% BEGIN O74E3
\paragraph{O74E3. The singular terms and the surviving coefficient.}
Suppose $4r<n$, $M>0$ and an actual level system
$(\Delta,\mathcal L,\Phi,\mu)$ is given in the sense of
\texttt{IsSystem}. Fix $L,L'\in\mathcal L$, an admissible lattice
$\Lambda\in V(E,n,M)$, and frames $x,y$ in the domain of $X$.
Put $\phi=\Phi(L,\Lambda)$, $\psi=\Phi(L',\Lambda)$.
For independent elliptic variables $(z,w)$ let $p$ denote the
doubled variables $(\mathrm{dblZ},\mathrm{dblz},\mathrm{dblz}')$
at $(x,y;z,w)$ and $k=k_0(x,y;z,w)$ the native Gaussian prefactor.
Use the source Weil and theta transformation laws and the graph
coordinate change. The contribution of $U$ to $k\,\mathrm{Eis}(D(\phi,\psi);p)$
is zero unless $U$ is the graph of an element $\gamma\in\Delta$.
For such a graph its entire contribution is
\[
 a\,\langle\Phi(L,\Lambda),\Phi(\gamma L',\Lambda)\rangle
       \,\overline{\mu(\gamma,L')}\; k_0(x,\gamma y;z,w).
\]

For the singular case of O74E1 use the system's cusp condition on
$\xi$ and $\phi$. It supplies a finite combination
$\sum_v b_v H_G(v\xi^t)\phi=0$ with $\sum_v b_v\ne0$.
Write $U$ as the bottom row space of $\sigma$. Every row of
$(v\xi^t,0)T$ belongs to $U$, so the first block of
$(v\xi^t,0)T\sigma^{-1}$ is zero. The Weil covariance and the
doubling law therefore replace each displacement by pure modulation
after applying $\omega(\sigma)$ to the doubled function. Evaluation
at zero is unaffected. Linearity in the first test gives
$(\sum_v b_v)(\omega(\sigma)D(\phi,\psi))(0)=0$.
This kills the numerator of the Eisenstein summand.

For a graph one may use $i_\gamma$ as representative. Indeed if
$\sigma,\sigma'$ have the same bottom space, $P=\sigma'\sigma^{-1}$
has lower left block zero. Write its blocks as $(A,B;0,D_0)$, so
$D_0=(A^*)^{-1}$. Its factor of automorphy is the constant
$d_0=\prod_{\tau\in\Psi}\tau(\det D_0)^n\ne0$. At the origin its
Weil action also multiplies by $d_0$. To check the latter formula,
the defining theta transformation for $P$ changes a summation matrix
$t$ into $tA$ and multiplies the test by $d_0$ and the quadratic
character with matrix $-BA^*$. In fact the new period is
$(AZ+B)A^*$, whose extra term $BA^*$ cancels that character, while
the linear elliptic terms are carried by the same change of index.
At $t=0$ the character equals one. The Weil composition law and the
cocycle of $\mathrm{hFac}$ thus multiply numerator and denominator
by the same $d_0$ under $\sigma\mapsto P\sigma$.

Write $p_\gamma,k_\gamma$
for the same variables and scalar with $y$ replaced by $\gamma y$,
and $j_\gamma=\mathrm{hFac}(i_\gamma;p)$. The native graph coordinate
law gives $i_\gamma p=p_\gamma$ and $k_\gamma j_\gamma=s_\gamma k$
for a nonzero constant $s_\gamma$. The defining theta product and the
ordinary unitary covariance $\Theta_{w_G(\gamma)\psi}(\gamma y)=\Theta_\psi(y)$
give
\[
 \Theta_H(\omega(i_\gamma)D(\phi,\psi);p_\gamma)
  =s_\gamma\Theta_H(D(\phi,w_G(\gamma)\psi);p_\gamma).
\]
Indeed the left side is $j_\gamma$ times the original doubled theta,
and multiplication by $k_\gamma$ uses precisely the displayed scalar
law. Varying $(z,w)$ gives all doubled elliptic variables. At fixed
period, the distinct rational matrix frequencies in $\Theta_H$ have
nonzero Gaussian factors and separate the coefficients by Fourier
uniqueness on a common period torus. Thus the tests themselves agree
up to $s_\gamma$. The graph summand is
$k(\omega(i_\gamma)D)(0)/j_\gamma
 =k_\gamma D(\phi,w_G(\gamma)\psi)(0)$.
By O74E2 the latter test value is $a\langle\phi,w_G(\gamma)\psi\rangle$.
The system's support law makes this zero for $\gamma\notin\Delta$.
Inside $\Delta$ the second test is $\mu(\gamma,L')\Phi(\gamma L',\Lambda)$.
Conjugate linearity in that slot gives the stated conjugation.
% END O74E3

% BEGIN O74E4
\paragraph{O74E4. The weighted sum and its elliptic derivatives.}
As convergence input it suffices to take absolute summability of the
native Eisenstein series at \emph{zero} elliptic variables in the
stable range. Indeed fix its period $Z$. At one place let $C,D_0$
be the bottom blocks of a representative, $A=CZ+D_0$, and
$B=A^{-1}C$. The unitary row relation $CD_0^*=D_0C^*$ gives
\[
 -\operatorname{Im}B=B(\operatorname{Im}Z)B^*.
\]
This follows by multiplying $B-B^*$ on the left by $A$ and on the
right by $A^*$. If $\operatorname{Im} Z\ge\delta I$, it gives
$\delta\|B\|^2\le\|\operatorname{Im}B\|\le\|B\|$ in operator norm.
A common $\delta>0$ works at the finitely many places, so for fixed
doubled elliptic variables the inverse Gaussian factor in each term
is uniformly bounded. Explicitly at one place the bound is
$\exp(2\pi\delta^{-1}\sum_k\|\zeta'_k\|\|\zeta_k\|)$.
Both the determinant and the finite test value are unchanged when
these variables are set to zero. The bound therefore transfers
absolute summability to every $(z,w)$. For $M>0$ use the finite class set of
the level $M$ and its chosen representatives in the definition of
$\operatorname{classSum}_M$. A representative has weight equal to the
reciprocal of the order of its stabilizer in the source group $D_H(M)$.
Set
\[
 c_\gamma=a\,\overline{\mu(\gamma,L')}
 \operatorname{classSum}_M\bigl(\Lambda\mapsto
  \langle\Phi(L,\Lambda),\Phi(\gamma L',\Lambda)\rangle\bigr).
\]
By O74E1 and O74E3, extending the graph terms by zero gives the
original summands for every doubled isotropic space. Absolute
summability allows restriction to this injectively indexed family.
The finite class sum commutes with it. Consequently
\[
 S(z,w):= k\,\operatorname{classSum}_M\bigl(\Lambda\mapsto
   \mathrm{Eis}(D(\Phi(L,\Lambda),\Phi(L',\Lambda));p)\bigr)
    =\sum_{\gamma\in\Delta} c_\gamma k_0(x,\gamma y;z,w).
\]
In particular the stabilizer weights belong to the definite lattice
classes. Each group element labels one graph, including when distinct
elements act identically on $L'$.

The series may be differentiated termwise to every finite order in
both elliptic variables. To see this, put $v=(z,w)$. The source
formula for $k_0$ has a constant determinant prefactor and the
exponential of a homogeneous complex quadratic $Q_\gamma(v)$, since
both doubled elliptic variables are linear in $v$. Write a term as
$A_\gamma\exp Q_\gamma(v)$. In the real dual of the finite-dimensional
space of complex homogeneous quadratics choose a basis from the
functionals $e_u:Q\mapsto\Re Q(u)$. They span the dual, since their
common kernel is zero (use $e^{\pi i/4}u$ to detect the imaginary
part). Their negatives are $e_{iu}$. For compact $K$, the coordinates
of $e_v$ in that basis have uniformly bounded sum of absolute values
as $v$ ranges over $K$. Scaling all basis vectors by a common constant
therefore puts $e(K)$ inside the convex hull of finitely many $e_{u_j}$:
a positive multiple of $e_u$ equals $e_{\sqrt{t}u}$ for multiplier $t$.
In zero dimension use the single probe $u_1=0$. It follows that
\[
 \sum_\gamma\sup_{v\in K}|A_\gamma e^{Q_\gamma(v)}|
 \le \sum_j\sum_\gamma |A_\gamma e^{Q_\gamma(u_j)}|<\infty .
\]
Cauchy's formula on a larger polydisc gives the same convergence for
each prescribed derivative on a smaller one. This proves the claim.

In particular fix $p_0+q_0=r$, row weights $c,c'$, and a finite
family $(P_j)$ of the corresponding row-highest harmonic polynomials
satisfying the source unitary column-invariance condition
\texttt{InvFam}. Let $\mathcal D_P$ be the sum over $j$ of
$P_j(\partial_z)\overline{P_j}(\partial_w)$ at zero in the original
letter directions. Since $K_P=\mathcal D_P k_0$, the differentiated
Eisenstein side is $\sum_\gamma c_\gamma K_P(x,\gamma y)$. To obtain
(ICL.6) one still needs to replace the Eisenstein side by the theta
side. Precisely, put
\[
 H(Z,\zeta,\zeta')=\operatorname{classSum}_M
   \bigl(\Lambda\mapsto(\Theta_H-\mathrm{Eis})
       (D(\Phi(L,\Lambda),\Phi(L',\Lambda));Z,\zeta,\zeta')\bigr).
\]
For the pair $L,L'$ the system's admission clause supplies $R$ for
this doubled family. O74E5 below proves the native heat equations for $H$. The
remaining two inputs here are the definite Siegel--Weil identity
$H(Z,0,0)=0$ for every allowed $Z$ and the invariant-polynomial
heat identity that these together imply $\mathcal D_P(kH(p))=0$ when $r\le n$. Under these
inputs the theta product formula identifies $S+kH(p)$ with the finite
weighted sum of
$\Theta_{\Phi(L,\Lambda)}(x;z)\overline{\Theta_{\Phi(L',\Lambda)}(y;\bar w)}$.
The real letter directions and conjugated polynomial in the second
slot give exactly the weighted sum of the differentiated theta
products in (ICL.6). The scalar $a$ was fixed in O74E2 before choosing
any level system, polynomials, lattices, or frames. Existence of the
level system, the stated Weil, theta and convergence inputs, and the
two remaining assertions about $H$ retain their separate hypotheses.
% END O74E4

% BEGIN O74E5
\paragraph{O74E5. The heat equations for the actual weighted difference.}
The function $H$ of O74E4 is holomorphic on the open product of the
Hermitian half spaces and the two elliptic matrix spaces, and satisfies,
at each embedding $\tau$ and for every $a,b$,
\[
 \lambda\frac{\partial H}{\partial Z_{ab}}
  =\sum_{k=1}^n\frac{\partial^2 H}
       {\partial\zeta_{ka}\,\partial\zeta'_{kb}},\qquad \lambda=2\pi i.
\]
All derivatives are complex coordinate derivatives at $\tau$ with
the other embeddings fixed. The two elliptic matrices are independent.
We use the zero-argument stable-range convergence input of O74E4 at
every allowed period. This suffices also for locally uniform convergence
in the period.

For the theta side write $u_k$ for the row of the summation matrix at
$\tau$. Its exponential there is
$\exp\{\lambda\sum_k(u_k Z\bar u_k^t+u_k\zeta_k+
 \zeta'_k\bar u_k^t)\}$, with $\zeta_k$ a column and $\zeta'_k$ a row.
Each side of the differential equation multiplies the summand by
$\lambda^2\sum_k u_{ka}\overline{u_{kb}}$.
To justify the derivatives, on a small closed polydisc each
$\operatorname{Im} Z_\tau$ has a common positive lower bound and both
elliptic matrices are bounded. The modulus of a term is then bounded
by a constant times $\exp(-c\|u\|^2+C\|u\|)$ with $c>0$, taking
all embeddings of the CM type together. The fixed finite test is bounded
and supported in a fixed denominator lattice in Minkowski space by
its Schwartz--Bruhat level. A Gaussian with this linear error is
summable over that lattice, giving local normal convergence.

For an Eisenstein summand keep the notation $A=CZ+D_0$ and
$B=A^{-1}C$ at the given place. Its variable factor is
\[
 F=(\det A)^{-n}\exp\left(-\lambda\sum_k\zeta'_k B\zeta_k\right).
\]
For an increment $V$ in $Z$, the differential of $B$ is $-BVB$,
and the logarithmic differential of $\det A$ is $\operatorname{tr}(BV)$.
Take $V$ to have its only nonzero entry $1$ in position $(a,b)$.
Dividing $\lambda\partial_{Z_{ab}}F$ by the nonzero factor $F$ gives
\[
 -n\lambda B_{ba}+\lambda^2\sum_k(\zeta'_k B)_a(B\zeta_k)_b .
\]
The mixed derivative for a single $k$, divided by $F$, contributes
$-\lambda B_{ba}+\lambda^2(\zeta'_kB)_a(B\zeta_k)_b$.
Summing gives the same expression. The test value at zero and all
other archimedean factors are constant for these derivatives.

Fix a central period $Z_*$ and put $B_*=(CZ_*+D_0)^{-1}C$ at each
place. If $\operatorname{Im} Z_*\ge\delta I$, O74E4 gives
$\|B_*\|\le\delta^{-1}$ uniformly in the representative. For
$Z=Z_*+V$ with $\|V\|<\delta/2$, factor
\[
 CZ+D_0=(CZ_*+D_0)(I+B_*V).
\]
All singular values of the second factor are at least $1/2$.
The inverse determinant to power $n$ grows by at most $2^{rn}$
at each place. Also $\operatorname{Im} Z\ge(\delta/2)I$; the
elliptic exponential has a uniform bound on every bounded set of
elliptic arguments by O74E4. The absolute terms at $(Z_*,0,0)$
therefore supply a single summable majorant, up to a constant,
throughout a neighborhood of any prescribed point. Cauchy's formula
now allows the holomorphic sums and their derivatives to be taken
termwise. The finite weighted difference over the definite lattice
classes preserves the equation and holomorphy. This gives precisely
\texttt{IsHeat} for the original $H$.
% END O74E5
